\documentclass[12pt]{article}
\usepackage[utf8]{inputenc}
\usepackage[T1]{fontenc}
\usepackage[margin=1in]{geometry}
\usepackage{setspace}
\usepackage{xcolor}
\usepackage[colorlinks=true, urlcolor=blue!60!black, linkcolor=blue!60!black]{hyperref}
\setlength{\parindent}{0em}
\setlength{\parskip}{0.6em}
\renewcommand{\baselinestretch}{1.0}
\begin{document}
\begin{center}
{\LARGE \textit{Everson v. Board of Education}, 330 U.S. 1 (1947)\par}
\vspace{0.5em}
{\large Case Brief --- Two-Page Edition\par}
\end{center}

\section{Memory Jogger}
The Court incorporated the Establishment Clause and proclaimed a wall between church and state, yet upheld 5--4 New Jersey's bus-fare reimbursements for parochial-school children as a neutral welfare benefit (\href{https://supreme.justia.com/cases/federal/us/330/1/}{330 U.S. 1}).

\section{Detailed Case Facts}
Under a 1941 New Jersey statute, Ewing Township reimbursed parents from tax funds for public-bus fares to public high schools and four Catholic parochial schools, which gave religious instruction alongside secular education. Taxpayer Arch R. Everson sued.

\section{Procedural History}
An intermediate state court held the payments unauthorized; the Court of Errors and Appeals reversed, and Everson appealed the federal questions.

\section{Judicial Votes}
Affirmed, 5--4: Black wrote for Vinson, Reed, Douglas, and Murphy; Jackson (with Frankfurter) and Rutledge (with Frankfurter, Jackson, and Burton) dissented.

\section{Holding}
The reimbursements serve a public purpose. The Establishment Clause binds the states through the Fourteenth Amendment, but a neutral benefit paid to parents of all faiths, like police and fire protection, does not breach the wall of separation.

\section{Analysis of Opinions}
Reading the Clause through Madison's \emph{Memorial and Remonstrance} and Jefferson's Virginia statute, Black listed strict prohibitions, including any tax to support religious institutions. But the state must be neutral, not religion's adversary, and cannot exclude believers from public welfare benefits; services ``so separate and so indisputably marked off from the religious function'' are permissible. Jackson found the separationist rhetoric ``utterly discordant with its conclusion'' and noted that only public and Catholic school students were reimbursed. Rutledge, appending Madison's \emph{Remonstrance}, argued that the Clause forbids ``every form of public aid or support for religion'' and that no rational line separates buses from books or salaries.

\section{Examples: Future Applications}
\textbf{Same-Side:} free transit passes for every student, including at religious schools, are valid. \textbf{Opposite-Side:} tax grants to all faiths to pay clergy salaries are invalid. \textbf{Fence-Sitter:} whether a State may require charter schools to be nonsectarian after \textit{Carson}.

\section{Critique}
Critics argue the original Clause protected state establishments from federal interference, making incorporation problematic, and that Virginia's struggle was not the sole gloss on 1789. Rutledge's fungibility point and Jackson's warning that government regulates what it subsidizes still frame aid disputes.

\section{Key Quotations}
\begin{quote}
``The `establishment of religion' clause of the First Amendment means at least this: neither a state nor the Federal Government can set up a church. Neither can pass laws which aid one religion, aid all religions, or prefer one religion over another. Neither can force nor influence a person to go to or to remain away from church against his will or force him to profess a belief or disbelief in any religion. No person can be punished for entertaining or professing religious beliefs or disbeliefs, for church attendance or non-attendance. No tax in any amount, large or small, can be levied to support any religious activities or institutions, whatever they may be called, or whatever form they may adopt to teach or practice religion.''
\end{quote}
\noindent --- Justice Black, Opinion of the Court, 330 U.S. at 15--16.

\end{document}
